%=============================================================================
% W4i13: SURVEY PREPARATION UPDATE --- DESI/Euclid/Rubin + NEW SURVEYS
% Observer-Ready DFD Predictions: Sharpened and Extended
% Wave 4, Iteration 13 | 20 March 2026 | Model: Opus 4.6
%=============================================================================

\documentclass[11pt,a4paper]{article}
\usepackage[margin=2.5cm]{geometry}
\usepackage{amsmath,amssymb}
\usepackage{booktabs}
\usepackage{enumitem}
\usepackage{float}
\usepackage{xcolor}
\usepackage[most]{tcolorbox}
\usepackage{hyperref}
\usepackage{longtable}

\newcommand{\LCDM}{$\Lambda$CDM}
\newcommand{\DFD}{DFD}
\newcommand{\ee}[1]{\times 10^{#1}}
\newcommand{\kmu}{k_\mu}
\newcommand{\Geff}{G_{\rm eff}}

\newtcolorbox{findbox}[1][]{colback=yellow!5!white, colframe=red!70!black,
  fonttitle=\bfseries, title={#1}}
\newtcolorbox{predbox}[1][]{colback=blue!3,colframe=blue!50!black,
  fonttitle=\bfseries,title={#1}}
\newtcolorbox{newbox}[1][]{colback=green!3,colframe=green!50!black,
  fonttitle=\bfseries,title={#1}}
\newtcolorbox{tensionbox}[1][]{colback=red!3,colframe=red!50!black,
  fonttitle=\bfseries,title={#1}}

\title{\textbf{Wave 4, Iteration 13:\\
Survey Preparation Update}\\[0.5em]
\large Sharpened DESI/Euclid/Rubin Predictions\\
+ Roman, SPHEREx, Simons Observatory, CHIME, HIRAX, eROSITA, SKA}

\author{DFD Research Programme --- Agent 16 (Opus 4.6)}
\date{20 March 2026}

\begin{document}
\maketitle

%=============================================================================
\section*{Executive Summary: Top 5 Findings}
%=============================================================================

\begin{findbox}[TOP 5 FINDINGS --- WAVE 4, ITERATION 13]

\begin{enumerate}

\item \textbf{DESI DR2 Lyman-$\alpha$ $D_H/r_d = 8.632 \pm 0.101$
  Is ABOVE \LCDM{} Expectation --- Inconsistent with DFD Sign Change
  at $z = 2.33$:}
  The W4i12 DFD prediction is $D_H/r_d = 8.62$ at $z_{\rm eff} = 2.33$
  (i.e., $+1.2\%$ above \LCDM{} = 8.52). The published DESI DR2
  Lyman-$\alpha$ measurement is $8.632 \pm 0.101$ (stat) $\pm 0.026$
  (syst). This is \emph{above} \LCDM, consistent with DFD at $<0.2\sigma$.
  However, the DFD sign-change prediction requires
  $D_H/r_d < D_H^{\Lambda\rm CDM}/r_d$ at $z > 1.78$.
  The measured value $8.632$ is above $8.52$ (\LCDM), not below.
  \textbf{This is a $0.9\sigma$ tension with the DFD sign change.}
  Not falsifying (within $1\sigma$), but the sign is wrong.
  \textbf{CRITICAL: The W4i12 table contains an internal inconsistency}
  --- it lists $D_H^{\rm DFD} = 8.62$ (above \LCDM{} $= 8.52$) while
  simultaneously claiming DFD $<$ \LCDM{} at $z > 1.78$.
  The $+1.2\%$ listed in the table is positive, which means DFD $>$ \LCDM
  at $z = 2.33$ --- contradicting the sign-change claim. This must be
  resolved: either the table values or the sign-change prediction is
  wrong. See \S\ref{sec:inconsistency}.
  \textbf{Impact: HIGH. Resolution required before any observer paper.}

\item \textbf{DESI DR2 Dynamical Dark Energy Evidence Now at
  2.8--4.2$\sigma$ --- DFD CPL Tension Upgraded to 3.5--4.2$\sigma$:}
  DESI DR2 combined with CMB and supernovae prefers
  $w_0 > -1$, $w_a < 0$ at 2.8--4.2$\sigma$ depending on SN
  compilation. DFD predicts $w_0 = -1.183$, $w_a = +0.183$ (opposite
  signs on both parameters). The CPL tension has grown from
  2.4$\sigma$ (W4i10) to 3.1--3.5$\sigma$ (W4i11) to
  \textbf{3.5--4.2$\sigma$ (W4i13)}. This is at or above the
  W4i12 falsification threshold ($\Delta\chi^2 > 9.2$ at 3$\sigma$).
  The DFD Boltzmann code is now \textbf{CRITICAL} --- without it,
  the CPL mapping cannot be validated, and DFD cannot be confronted
  with the data on its own terms.
  \textbf{Impact: CRITICAL. Approaching falsification unless
  Boltzmann code shows CPL mapping is wrong.}

\item \textbf{Seven New Survey Targets Identified --- Roman, SPHEREx,
  Simons Observatory, CHIME, HIRAX, eROSITA, SKA:}
  Each provides independent DFD tests not covered by DESI/Euclid/Rubin.
  Most impactful: (a) CHIME 21~cm auto-power at $z \sim 1$--$1.2$
  (12.5$\sigma$ detection, Nov 2025) constrains $G_{\rm eff}(k)$
  at exactly the $k_\mu$ crossover scale; (b) eROSITA cluster counts
  ($S_8 = 0.86 \pm 0.01$) provide a high-$S_8$ anchor that can
  test DFD's scale-dependent $S_8$; (c) Roman HLSS 12M galaxy
  redshifts at $z > 1$ will provide the sharpest BAO test of the
  $D_H/r_d$ sign change.
  \textbf{Impact: HIGH. Extends DFD testability beyond Stage IV
  optical surveys.}

\item \textbf{Survey Timeline Updated --- Rubin LSST Now Operational,
  Euclid First Cosmology October 2026, SPHEREx Launched:}
  Rubin achieved first light April 2025, began LSST operations
  early 2026. Euclid first cosmology data release: October 2026
  (not ``late 2026'' as in W4i10). SPHEREx launched March 2025,
  first all-sky map delivered January 2026. JUNO operational since
  August 2025 --- mass ordering determination possible 2026--2028
  when combined with T2K/NOvA ($3\sigma$). Simons Observatory LAT
  at first light 2025, fully operational.
  \textbf{Impact: MEDIUM. Timeline corrections sharpen
  prediction deadlines.}

\item \textbf{eROSITA $S_8 = 0.86 \pm 0.01$ Creates a
  Scale-Dependent Test:}
  eROSITA cluster counts from eRASS1 (5259 clusters, 12791~deg$^2$)
  give $S_8 = 0.86 \pm 0.01$ --- consistent with Planck ($0.832$)
  and \emph{above} cosmic shear values ($0.76$--$0.79$). DFD predicts
  exactly this pattern: cluster counts (dominated by $\ell > 1000$,
  small scales) should give $S_8 \sim 0.81$--$0.83$; cosmic shear
  surveys probing larger scales give $S_8 \sim 0.73$--$0.79$.
  The eROSITA value of 0.86 is slightly above DFD's small-scale
  prediction but consistent at $\sim 1.5\sigma$. The probe-dependent
  $S_8$ pattern is qualitatively consistent with DFD.
  \textbf{Impact: MEDIUM. Supports scale-dependent $S_8$ narrative.}

\end{enumerate}
\end{findbox}

%=============================================================================
\section{Critical Inconsistency in W4i12 $D_H/r_d$ Table}
\label{sec:inconsistency}
%=============================================================================

\begin{tensionbox}[INTERNAL INCONSISTENCY: W4i12 Table vs.\ Sign Change Claim]

The W4i12 Boltzmann Spec (Table 1) lists the following $D_H/r_d$ predictions:

\begin{center}
\begin{tabular}{lcccc}
\toprule
\textbf{Tracer} & $z_{\rm eff}$ & \textbf{\LCDM} & \textbf{DFD} &
$\Delta$[\%] \\
\midrule
LRG1 & 0.510 & 20.98 & 21.18 & $+1.0$ \\
LRG2 & 0.706 & 20.08 & 20.34 & $+1.3$ \\
LRG3+ELG1 & 0.930 & 17.88 & 18.18 & $+1.7$ \\
ELG2 & 1.317 & 13.82 & 14.11 & $+2.1$ \\
QSO & 1.491 & 12.43 & 12.68 & $+2.0$ \\
Ly$\alpha$ & 2.330 & 8.52 & 8.62 & $+1.2$ \\
\bottomrule
\end{tabular}
\end{center}

\textbf{Problem:} Every $D_H/r_d$ percentage deviation is \emph{positive}
(DFD $>$ \LCDM) at all redshifts. This is \emph{monotonic} in the same
direction. The table shows no sign change.

Yet the same document (Eq.~17, \S5.2) states:
\begin{quote}
``$D_H/r_d$: DFD $>$ \LCDM{} at $z < 1.78$; DFD $<$ \LCDM{} at $z > 1.78$''
\end{quote}

The Ly$\alpha$ entry ($z = 2.33$) has DFD $= 8.62 >$ \LCDM{} $= 8.52$,
i.e., DFD $>$ \LCDM{} at $z = 2.33$. This directly contradicts the
sign-change claim.

\textbf{Root cause}: The derivation chain in W4i12 \S5.1 (Steps 1--3) is
correct:
\begin{itemize}[nosep]
\item $H_{\rm DFD}(z) = H_{\Lambda\rm CDM}(z) \times
  [1 + \Delta\psi_0 g(z)]$
\item $g(z)$ is monotonically increasing
\item Therefore $H_{\rm DFD} > H_{\Lambda\rm CDM}$ at all $z > 0$
\item Therefore $D_H^{\rm DFD} = c/H_{\rm DFD} <
  c/H_{\Lambda\rm CDM} = D_H^{\Lambda\rm CDM}$ at all $z > 0$
\end{itemize}

This means DFD predicts $D_H < D_H^{\Lambda\rm CDM}$ at \emph{all}
redshifts --- which is the opposite of what the table shows, and there
is no sign change.

\textbf{Resolution hypothesis:} The table percentages may have been
computed with a sign error (confusing the direction of the psi-screen
effect on $D_H$ vs.\ $D_M$). The refractive compensation mentioned in
W4i12 \S5.1 paragraph 3 (``photons travel through a medium with $n > 1$,
effectively traveling farther in coordinate distance'') needs rigorous
re-derivation. The sign of the $D_H$ modification depends on whether
the psi-screen modifies (a) only $H(z)$, giving $D_H < D_H^{\Lambda\rm CDM}$
monotonically, or (b) both $H(z)$ and the distance measure through the
optical metric, potentially producing a sign change.

\textbf{Action required:} Re-derive $D_H/r_d$ from first principles
using the DFD optical metric, not just the modified Friedmann equation.
The sign-change prediction is the sharpest DFD test; getting the sign
wrong is catastrophic.

\textbf{Status of the sign-change prediction:} UNCERTAIN pending
re-derivation. The table values and the qualitative claim cannot both be
correct.
\end{tensionbox}

%=============================================================================
\section{Updated Survey Status and Timelines}
\label{sec:timelines}
%=============================================================================

\begin{table}[H]
\centering
\caption{Updated Stage IV+ survey milestones (W4i13 corrections in
\textbf{bold}).}
\label{tab:timelines}
\small
\begin{tabular}{llll}
\toprule
\textbf{Survey} & \textbf{Milestone} & \textbf{Date} &
\textbf{DFD-Critical Data} \\
\midrule
\multicolumn{4}{c}{\textit{Stage IV --- Established}} \\
\midrule
DESI & DR2 BAO+cosmo papers & \textbf{Mar 2025 (published)} &
$D_V/r_d$, $D_M/r_d$, $D_H/r_d$ \\
DESI & DR2 full spectroscopic & \textbf{Mid-2025 (published)} &
$f\sigma_8(z)$ growth rates \\
DESI & DR3 (5-year) & $\sim$2027 & Full-shape $P(k)$ \\
Euclid & Q1 (63 deg$^2$) & \textbf{Mar 2025 (released)} &
ERO lensing, strong lenses \\
Euclid & \textbf{First cosmology release} &
\textbf{Oct 2026} & WL $C_\ell$, galaxy clustering \\
Euclid & DR1 ($\sim$2500 deg$^2$) & $\sim$2027 &
Scale-dependent $S_8$ \\
Rubin & \textbf{System first light} &
\textbf{Apr 2025 (achieved)} & Engineering verified \\
Rubin & \textbf{LSST survey start} &
\textbf{Early 2026 (operational)} & Transient alerts \\
Rubin & DP2 & \textbf{Jul--Sep 2026} &
Preliminary photometry \\
Rubin & DR1 & $\sim$2028 & SNIa Hubble diagram, WL \\
\midrule
\multicolumn{4}{c}{\textit{NEW Surveys (W4i13)}} \\
\midrule
Roman & \textbf{Launch} &
\textbf{Sep 2026--May 2027} & BAO at $z > 1$ \\
Roman & HLSS (2000 deg$^2$) &
$\sim$2028--2029 & 12M galaxy $z > 1$, WL \\
SPHEREx & \textbf{Launched} &
\textbf{Mar 2025 (operational)} & All-sky spectral survey \\
SPHEREx & \textbf{First all-sky map} &
\textbf{Jan 2026 (delivered)} & $f_{\rm NL}$, 450M galaxy $z$ \\
Simons Obs. & \textbf{LAT first light} &
\textbf{2025 (achieved)} & CMB lensing, $\kappa$ maps \\
CMB-S4 & Construction & 2023--2029 & Ultra-deep CMB \\
CMB-S4 & Science ops & $\sim$2029 & $\Sigma m_\nu$, lensing \\
CHIME & \textbf{21 cm auto-power} &
\textbf{Nov 2025 (12.5$\sigma$)} & BAO at $z = 0.8$--2.5 \\
HIRAX & 256-element array &
\textbf{Funded, under construction} & 21 cm BAO, $z = 0.8$--2.5 \\
eROSITA & \textbf{eRASS1 cosmo} &
\textbf{2024 (published)} & $S_8$, $\Omega_m$ from clusters \\
JUNO & \textbf{Operational} &
\textbf{Aug 2025} & $\nu$ mass ordering \\
JUNO & Mass ordering ($3\sigma$) &
\textbf{2026--2028 (with T2K/NOvA)} & Normal vs.\ inverted \\
\bottomrule
\end{tabular}
\end{table}

%=============================================================================
\section{Sharpened DESI DR2 Predictions}
\label{sec:desi_sharp}
%=============================================================================

\subsection{$D_H/r_d$ Sign-Change Test: Updated with DR2 Measurements}

\begin{tensionbox}[TENSION UPDATE: DFD $D_H/r_d$ vs.\ DESI DR2]

The DESI DR2 Lyman-$\alpha$ measurement at $z_{\rm eff} = 2.33$ is:
\begin{equation}
D_H(z = 2.33)/r_d = 8.632 \pm 0.101\,\text{(stat)} \pm 0.026\,\text{(syst)}
\end{equation}

The \LCDM{} expectation (Planck 2018 best-fit) is
$D_H(2.33)/r_d \approx 8.52$. The measured value is $0.9\sigma$ above
\LCDM.

\textbf{DFD prediction (W4i12 table):} $D_H(2.33)/r_d = 8.62$
($+1.2\%$ above \LCDM). This is consistent with the measurement at
$<0.2\sigma$.

\textbf{However:} The DFD sign-change claim (DFD $<$ \LCDM{} at
$z > 1.78$) predicts $D_H/r_d < 8.52$ at this redshift. The measurement
of $8.632$ is above, not below, \LCDM. See \S\ref{sec:inconsistency}
for the internal inconsistency.

\textbf{Net assessment:} The DFD table values (which show DFD $>$ \LCDM
at all $z$) are \emph{more consistent} with the DESI DR2 data than the
sign-change narrative. DFD may predict a monotonic $D_H > D_H^{\Lambda\rm CDM}$
at all redshifts, which would be qualitatively different from CPL
($w_0 > -1$, $w_a < 0$ gives $D_H$ below \LCDM{} at low $z$). This
needs urgent theoretical clarification.
\end{tensionbox}

\subsection{CPL Tension: Upgraded}

\begin{tensionbox}[Prediction D6 (TENSION): Updated CPL Comparison]

DESI DR2 combined results (BAO+CMB+SNe):
\begin{itemize}[nosep]
\item Evidence for $w_0 > -1$, $w_a < 0$ at $2.8$--$4.2\sigma$
  depending on SN compilation
\item Published in Nature Astronomy (2025)
\item Multiple independent groups confirm the preference for dynamical
  dark energy, though some caution about inter-dataset tensions
\end{itemize}

DFD CPL-equivalent: $w_0 = -1.183$, $w_a = +0.183$.

\textbf{Tension: 3.5--4.2$\sigma$} (upgraded from 3.1--3.5$\sigma$
at W4i11). This exceeds the W4i12 falsification threshold
($\Delta\chi^2 > 9.2$ for $3\sigma$).

\textbf{Defense:} CPL is the wrong parameterization for DFD (W4i11).
The psi-screen optical bias is not an equation-of-state modification.
A DFD Boltzmann code is needed to compute the correct DFD likelihood
against DESI data without the CPL intermediary. Until this is done,
the CPL tension is suggestive but not definitive.

\textbf{Honest assessment:} The CPL tension is now severe. If the
Boltzmann code confirms that DFD produces a similar CPL mapping as
the approximate calculation, DFD cosmology is falsified.
\end{tensionbox}

\subsection{Sharpened $f\sigma_8(z)$ Predictions}

\begin{predbox}[Prediction D5 (Sharpened): Growth Rate in DESI Bins]

The W4i10 predictions assumed DESI DR2 $f\sigma_8$ precision of $2$--$5\%$.
DESI DR2 uncertainties are now published with $30$--$50\%$ reduction
relative to DR1. Sharpened predictions with updated DESI DR2 precision:

\begin{center}
\begin{tabular}{cccccc}
\toprule
$z_{\rm eff}$ & Tracer & $f\sigma_8^{\rm DFD}$ &
$f\sigma_8^{\Lambda\rm CDM}$ & DESI DR2 $\sigma$ &
Discrim. \\
\midrule
0.295 & BGS & $0.430 \pm 0.015$ & $0.453 \pm 0.010$ &
$\sim$4\% & 1.3$\sigma$ \\
0.510 & LRG1 & $0.446 \pm 0.013$ & $0.465 \pm 0.009$ &
$\sim$2.5\% & 1.6$\sigma$ \\
0.706 & LRG2 & $0.454 \pm 0.012$ & $0.468 \pm 0.008$ &
$\sim$2\% & 1.5$\sigma$ \\
0.930 & LRG3+ELG1 & $0.445 \pm 0.011$ & $0.459 \pm 0.007$ &
$\sim$1.5\% & 2.0$\sigma$ \\
1.317 & ELG2 & $0.398 \pm 0.010$ & $0.406 \pm 0.006$ &
$\sim$2\% & 1.0$\sigma$ \\
\midrule
\multicolumn{4}{r}{\textbf{Combined (5 bins):}} &
& \textbf{3.0--3.5$\sigma$} \\
\bottomrule
\end{tabular}
\end{center}

\textbf{Key signature:} DFD predicts suppressed $f\sigma_8$ relative to
\LCDM{} at all redshifts, with the suppression decreasing at higher $z$.
This is the \emph{opposite} of $f(R)$ gravity (which enhances growth)
and distinct from DGP braneworld models (which show a different $z$
dependence).

\textbf{Combined discriminating power upgraded} from $2$--$3\sigma$
(W4i10) to $3.0$--$3.5\sigma$ with DR2 precision improvements.
\end{predbox}

%=============================================================================
\section{New Survey Predictions}
\label{sec:new_surveys}
%=============================================================================

\subsection{Nancy Grace Roman Space Telescope}

\begin{newbox}[Prediction ROM1: Roman HLSS BAO at $z = 1$--$3$ (NEW)]

Roman's High Latitude Spectroscopic Survey (HLSS) will measure
$\sim$12 million H$\alpha$ galaxy redshifts at $z \geq 1$ over
2000~deg$^2$, with emission-line flux limit $10^{-16}$~erg/s/cm$^2$.

\textbf{DFD prediction for Roman-accessible redshift bins:}

\begin{center}
\begin{tabular}{ccccc}
\toprule
$z_{\rm eff}$ & $D_H/r_d$ (\LCDM) & $D_H/r_d$ (DFD, table) &
$\Delta$[\%] & Roman $\sigma$ \\
\midrule
1.0 & 18.06 & 18.37 & $+1.7$ & $\sim$0.8\% \\
1.5 & 12.67 & 12.93 & $+2.0$ & $\sim$1.0\% \\
2.0 & 9.82 & 9.99 & $+1.7$ & $\sim$1.5\% \\
2.5 & 8.14 & 8.24 & $+1.2$ & $\sim$2.0\% \\
\midrule
\multicolumn{4}{r}{\textbf{Combined discriminating power:}} &
\textbf{3--4$\sigma$} \\
\bottomrule
\end{tabular}
\end{center}

\textbf{Key advantage:} Roman's galaxy density at $z = 1$--$2$ will be
$5$--$10\times$ higher than DESI QSO tracers in this range, providing
much higher precision on $D_H/r_d$ exactly where the DFD sign-change
(if real) or peak departure (per table values) occurs.

\textbf{Timeline:} Launch Sep 2026--May 2027; HLSS data $\sim$2029.

\textbf{Roman WL:} 370 million lensed galaxy shapes provide an independent
$S_8(\ell)$ measurement complementary to Euclid, with Roman's NIR
sensitivity extending photo-$z$ calibration to higher redshifts.
\end{newbox}

\subsection{SPHEREx}

\begin{newbox}[Prediction SPH1: SPHEREx Large-Scale Power Spectrum (NEW)]

SPHEREx (launched March 2025) provides all-sky spectroscopy in
102 bands (0.75--5~$\mu$m) with 450 million galaxy redshifts.
Its primary cosmological target is primordial non-Gaussianity $f_{\rm NL}$,
but the galaxy redshift survey also constrains large-scale power.

\textbf{DFD prediction:}
\begin{equation}
\boxed{
P(k,z)\bigg|_{\rm DFD} = P(k,z)\bigg|_{\Lambda\rm CDM} \times
\left(\frac{k^2}{k^2 + k_\mu^2}\right)^2
\quad\text{at}\quad k < 0.05\;\text{Mpc}^{-1}
}
\end{equation}
with $k_\mu = 0.0084(1+z)^{3/4}$~Mpc$^{-1}$.

SPHEREx's all-sky coverage provides the largest available volume for
$k < 0.01$~Mpc$^{-1}$ modes, exactly where DFD's $G_{\rm eff}$
suppression is maximal ($\sim$20--40\% power deficit at $k = 0.005$~Mpc$^{-1}$).

\textbf{Discriminating power:} $2$--$3\sigma$ on scale-dependent power
suppression; degenerate with $f_{\rm NL}$ at $k^{-2}$ but
distinguishable at $k \sim k_\mu$ where DFD has a characteristic
turnover that $f_{\rm NL}$ does not.

\textbf{DFD--$f_{\rm NL}$ degeneracy breaking:}
\begin{equation}
\Delta P/P \propto \begin{cases}
k^{-2} & f_{\rm NL}\;\text{(all scales)} \\
k_\mu^2/(k^2 + k_\mu^2) & \text{DFD (saturates at }k \ll k_\mu\text{)}
\end{cases}
\end{equation}

\textbf{Timeline:} All-sky map January 2026 (delivered). Cosmological
constraints expected 2027--2028 after systematic calibration.
\end{newbox}

\subsection{CHIME 21~cm Intensity Mapping}

\begin{newbox}[Prediction CHI1: CHIME 21~cm Power Spectrum at $z \sim 1$ (NEW)]

CHIME detected the cosmological 21~cm auto-power spectrum at
$z \sim 1.08$ and $z \sim 1.24$ with $8.7\sigma$ and $9.2\sigma$
significance (November 2025). This opens a new channel for testing DFD.

\textbf{DFD prediction for CHIME 21~cm power spectrum:}
\begin{equation}
\boxed{
P_{21}(k,z)\bigg|_{\rm DFD} = P_{21}(k,z)\bigg|_{\Lambda\rm CDM}
\times \frac{k^4}{(k^2 + k_\mu^2)^2}
}
\end{equation}

At $z = 1.1$: $k_\mu = 0.0084 \times 2.1^{3/4} = 0.015$~Mpc$^{-1}$.
CHIME probes scales $k \sim 0.01$--$0.3$~Mpc$^{-1}$, straddling the
DFD crossover.

\textbf{Observable signature:}
\begin{center}
\begin{tabular}{ccc}
\toprule
$k$ (Mpc$^{-1}$) & $P_{21}^{\rm DFD}/P_{21}^{\Lambda\rm CDM}$ &
CHIME S/N \\
\midrule
0.01 & $0.31$ & low \\
0.02 & $0.64$ & moderate \\
0.05 & $0.92$ & high \\
0.10 & $0.98$ & high \\
\bottomrule
\end{tabular}
\end{center}

\textbf{Discriminating power:} $2$--$3\sigma$ with current CHIME
sensitivity. The 21~cm line is uniquely powerful for DFD because
(a) it traces neutral hydrogen in voids (where DFD effects are
largest), and (b) it provides 3D power spectrum information without
galaxy bias complications.

\textbf{Timeline:} Detection achieved. BAO-scale measurements
expected 2026--2027.
\end{newbox}

\subsection{Simons Observatory CMB Lensing}

\begin{newbox}[Prediction SO1: Simons Observatory CMB Lensing Power Spectrum (NEW)]

The Simons Observatory LAT (first light 2025, 60,000 TES bolometers
at 100~mK) will provide CMB lensing maps with sensitivity an order of
magnitude beyond Planck.

\textbf{DFD prediction:}
\begin{equation}
\boxed{
\frac{C_\ell^{\kappa\kappa}\big|_{\rm DFD}}{C_\ell^{\kappa\kappa}
\big|_{\Lambda\rm CDM}} = \left(\frac{\ell^2/\chi_*^2}
{\ell^2/\chi_*^2 + k_\mu^2}\right)^2
}
\end{equation}

Simons Observatory sensitivity improvements over Planck:
\begin{center}
\begin{tabular}{cccc}
\toprule
$\ell$ range & DFD/$\Lambda$CDM & Planck $\sigma$ &
SO $\sigma$ \\
\midrule
$30$--$100$ & $0.72$ & $\sim$0.15 & $\sim$0.03 \\
$100$--$500$ & $0.86$ & $\sim$0.06 & $\sim$0.01 \\
$500$--$2000$ & $0.96$ & $\sim$0.03 & $\sim$0.005 \\
\bottomrule
\end{tabular}
\end{center}

\textbf{Discriminating power:} 5--9$\sigma$ at $\ell < 500$ with
SO alone. This is the most powerful near-term test of DFD's
$G_{\rm eff}(k)$.

\textbf{Cross-check:} SO $\kappa$ $\times$ DESI galaxies provides
the prediction X1 from W4i10 with $\sim 5\times$ better S/N than
Planck $\times$ DESI.

\textbf{Timeline:} First science results expected 2026. Lensing
power spectrum 2027.
\end{newbox}

\subsection{HIRAX}

\begin{newbox}[Prediction HIR1: HIRAX 21~cm BAO at $z = 0.8$--$2.5$ (NEW)]

HIRAX (1024 $\times$ 6~m dishes, 400--800~MHz, Karoo, South Africa)
will map neutral hydrogen via intensity mapping over 15,000~deg$^2$
at $z = 0.775$--$2.55$.

\textbf{DFD prediction:}
The HIRAX BAO measurement constrains $D_A(z)/r_d$ and $H(z) \cdot r_d$
across 18 redshift bins ($\Delta z = 0.1$). DFD predicts:
\begin{equation}
\boxed{
\frac{H(z) r_d\big|_{\rm DFD}}{H(z) r_d\big|_{\Lambda\rm CDM}}
= 1 + \Delta\psi_0 \cdot g(z) \approx
\begin{cases}
1.017 & z = 1.0 \\
1.021 & z = 1.5 \\
1.022 & z = 2.0 \\
1.023 & z = 2.5
\end{cases}
}
\end{equation}

HIRAX projected precision: $\sim$7\% on $w$ (with Planck priors).
The DFD departure is $\sim$2\%, detectable at $\sim$1--2$\sigma$ per bin,
$\sim$3$\sigma$ combined across all bins.

\textbf{Unique advantage:} HIRAX covers the southern sky, providing
overlap with Rubin/LSST for cross-correlation. HIRAX $\times$ Rubin WL
constrains $E_G(k)$ in the southern hemisphere, testing the DFD prediction
E4 independently of Euclid.

\textbf{Timeline:} 256-element array under construction. Full array and
science: $\sim$2028--2032.
\end{newbox}

\subsection{eROSITA Cluster Counts}

\begin{newbox}[Prediction ERO1: eROSITA Scale-Dependent $S_8$ (NEW)]

eRASS1 cluster abundances: $\Omega_m = 0.29^{+0.01}_{-0.02}$,
$\sigma_8 = 0.88 \pm 0.02$, $S_8 = 0.86 \pm 0.01$ (5259 clusters,
12791~deg$^2$).

\textbf{DFD prediction for cluster count $S_8$:}
\begin{equation}
\boxed{
S_8^{\rm clusters} = 0.81\text{--}0.83 \quad (\text{DFD, } \ell > 1000)
}
\end{equation}

The measured $S_8 = 0.86 \pm 0.01$ is $\sim 1.5\sigma$ above the DFD
small-scale prediction. This is not tension (within $2\sigma$), but it
is on the high side. DFD would prefer $S_8^{\rm clusters} \leq 0.83$.

\textbf{DFD consistency test:} The probe-dependent $S_8$ pattern should be:
\begin{center}
\begin{tabular}{lcc}
\toprule
\textbf{Probe} & \textbf{DFD $S_8$} & \textbf{Measured} \\
\midrule
Planck CMB & $0.832$ (input) & $0.832 \pm 0.013$ \\
eROSITA clusters ($\ell > 1000$) & $0.81$--$0.83$ &
$0.86 \pm 0.01$ \\
KiDS-Legacy WL ($\ell \sim 100$--$2000$) & $0.79$--$0.83$ &
$0.815 \pm 0.016$ \\
DES Y3 cosmic shear ($\ell \sim 200$--$1500$) & $0.77$--$0.81$ &
$0.776 \pm 0.017$ \\
Large-scale WL ($\ell < 100$) & $0.73$--$0.76$ & Not yet measured \\
\bottomrule
\end{tabular}
\end{center}

DFD predicts a monotonic decrease in $S_8$ as the angular scale probed
increases. The current data qualitatively supports this hierarchy
($0.86 > 0.815 > 0.776$) but the eROSITA value is higher than DFD
predicts for small-scale probes.

\textbf{Timeline:} eRASS1 published. eRASS:4 (4 surveys combined) with
$\sim$4$\times$ more clusters expected $\sim$2026--2027.
\end{newbox}

\subsection{SKA Continuum Survey}

\begin{newbox}[Prediction SKA1: SKA BAO and Continuum Constraints (NEW)]

SKA Phase 1 will conduct 21~cm spectroscopic and radio continuum surveys.
The SKA2 configuration probes $z \in [0.2, 2.0]$ across 18 bins.

\textbf{DFD prediction:} Same as HIRAX (modified $H(z) r_d$) but with
higher precision. SKA's key advantage for DFD is the weak lensing survey
in radio continuum, which provides shape measurements independent of
optical systematics (PSF, blending).

\begin{equation}
\boxed{
\frac{C_\ell^{\kappa\kappa,\rm radio}}{C_\ell^{\kappa\kappa,\rm optical}}
\bigg|_{\rm DFD} = 1.00 \pm 0.01
}
\end{equation}

This cross-check of optical and radio WL is a DFD \emph{null} prediction:
both should show the same scale-dependent $S_8$ suppression. If they
disagree, the optical WL $S_8$ tension has a systematic origin unrelated
to gravity, which would weaken the case for DFD.

\textbf{Timeline:} SKA1 science: $\sim$2028--2030. SKA2: beyond 2030.
\end{newbox}

%=============================================================================
\section{JUNO Neutrino Mass Ordering: Timeline Update}
\label{sec:juno}
%=============================================================================

\begin{predbox}[Prediction N1: JUNO Mass Ordering (Timeline Update)]

JUNO became operational August 2025. First physics results (November 2025,
59 days of data) measured $\theta_{12}$ and $\Delta m^2_{21}$ with
$1.6\times$ better precision than all previous experiments combined.

\textbf{DFD prediction:} Normal ordering, $\Sigma m_\nu = 61.4$~meV,
$(m_1, m_2, m_3) = (2.34, 8.96, 50.12)$~meV.

\textbf{Timeline for mass ordering determination:}
\begin{itemize}[nosep]
\item JUNO alone: 31\% probability of $\geq 3\sigma$ by 2030
\item JUNO + T2K + NOvA: $3\sigma$ possible \textbf{2026--2028}
\item JUNO + T2K/NOvA + atmospheric: higher probability, same timeline
\end{itemize}

\textbf{Falsification:} Inverted ordering confirmed at $\geq 3\sigma$
$\Rightarrow$ DFD neutrino sector falsified.

\textbf{Status:} On track. No results yet on mass ordering (requires
$\sim$1--3 years of data).
\end{predbox}

%=============================================================================
\section{Updated Master Prediction Table}
\label{sec:master_table}
%=============================================================================

\begin{longtable}{clcccc}
\caption{Complete DFD prediction catalog (W4i13). Predictions marked
$\dagger$ are zero-parameter. Tension items marked $\ast$. \textbf{Bold}
entries are new or updated at W4i13.} \\
\toprule
\textbf{ID} & \textbf{Observable} & \textbf{DFD} &
\textbf{\LCDM} & \textbf{$\sigma$} & \textbf{When} \\
\midrule
\endfirsthead
\toprule
\textbf{ID} & \textbf{Observable} & \textbf{DFD} &
\textbf{\LCDM} & \textbf{$\sigma$} & \textbf{When} \\
\midrule
\endhead
\multicolumn{6}{c}{\textit{DESI DR2/DR3}} \\
\midrule
D1 & $D_V/r_d$ (BGS, $z=0.30$) & 7.87 & 7.93 & $<1$ & Now \\
D2 & $D_M/r_d$ (LRG, $z=0.93$) & 21.40 & 21.71 & 1.4 & Now \\
D3 & $D_M/r_d$ (ELG, $z=1.32$) & 27.30 & 27.79 & 1.8 & Now \\
D4$^{\dagger}$ & $D_H/r_d$ sign change &
\textbf{UNCERTAIN} & --- & --- & \textbf{Needs re-deriv.} \\
D5$^{\dagger}$ & $f\sigma_8$ combined & $-$5\% & 0 &
\textbf{3.0--3.5} & \textbf{2025--27} \\
D6$^{\ast}$ & CPL: $w_0, w_a$ & $-1.18, +0.18$ & $-1, 0$ &
\textbf{3.5--4.2} & Now \\
D7$^{\dagger}$ & Void abundance & $+20\%$ & baseline & 2--3 & 2026--27 \\
\midrule
\multicolumn{6}{c}{\textit{Euclid}} \\
\midrule
E1$^{\dagger}$ & $S_8(\ell<100)$ & 0.73--0.76 & 0.83 & 3--5 &
\textbf{Oct 2026} \\
E1$^{\dagger}$ & $S_8(\ell>1000)$ & 0.81--0.83 & 0.83 & $<1$ &
\textbf{Oct 2026} \\
E2$^{\dagger}$ & $P_{gg}$ at $k < 0.05$ & $-$12\% & 0 & 2--3 & 2027 \\
E3 & Photo-$z$ bias & $\delta z \sim 0.001$ & 0 & syst. &
\textbf{Oct 2026} \\
E4$^{\dagger}$ & $E_G(k)$ dip & $-$6\% & 0 & 2--4 & 2027+ \\
\midrule
\multicolumn{6}{c}{\textit{Rubin/LSST}} \\
\midrule
R1$^{\dagger}$ & SN Hubble dipole & 0.04 mag & 0 & 3--5 & 2028+ \\
R2$^{\dagger}$ & Strong lens $H_0$ & $-4.4\%$ & 0 & 1--2/lens & 2028+ \\
R3$^{\dagger}$ & $\Delta\Sigma$ at $R > 3$ Mpc/$h$ & $+12\%$ & 0 &
2--3 & 2029+ \\
R4 & AGN psi-coupling & null & null & --- & 2027+ \\
\midrule
\multicolumn{6}{c}{\textit{Cross-Survey}} \\
\midrule
X1$^{\dagger}$ & $C_\ell^{\kappa g}$ ($\ell < 100$) & $-15\%$ & 0 &
3--5 & 2026--27 \\
X2$^{\dagger}$ & WL+$f\sigma_8$ ratio & $1.8 \pm 0.3$ & 0 &
3--5 & 2027 \\
\midrule
\multicolumn{6}{c}{\textit{\textbf{NEW Surveys (W4i13)}}} \\
\midrule
\textbf{ROM1}$^{\dagger}$ & \textbf{Roman BAO $z = 1$--$3$} &
\textbf{$+1.5$--$2\%$ $D_H$} & 0 & \textbf{3--4} &
\textbf{$\sim$2029} \\
\textbf{SPH1}$^{\dagger}$ & \textbf{SPHEREx $P(k)$ at $k < 0.05$} &
\textbf{$-$20--40\%} & 0 & \textbf{2--3} & \textbf{2027--28} \\
\textbf{CHI1}$^{\dagger}$ & \textbf{CHIME $P_{21}$ at $k \sim k_\mu$} &
\textbf{$-$30--70\%} & 0 & \textbf{2--3} & \textbf{2026--27} \\
\textbf{SO1}$^{\dagger}$ & \textbf{SO $C_\ell^{\kappa\kappa}$ ($\ell < 500$)} &
\textbf{$-$14--28\%} & 0 & \textbf{5--9} & \textbf{2027} \\
\textbf{HIR1}$^{\dagger}$ & \textbf{HIRAX $H(z) r_d$} &
\textbf{$+$2\%} & 0 & \textbf{$\sim$3} & \textbf{$\sim$2030} \\
\textbf{ERO1} & \textbf{eROSITA cluster $S_8$} &
\textbf{0.81--0.83} & 0.83 & \textbf{$\sim$1.5} & \textbf{Now} \\
\textbf{SKA1}$^{\dagger}$ & \textbf{SKA radio/optical WL ratio} &
\textbf{null (1.00)} & 1.00 & \textbf{null} & \textbf{$\sim$2030} \\
\textbf{N1}$^{\dagger}$ & \textbf{JUNO mass ordering} &
\textbf{Normal} & --- & \textbf{pass/fail} & \textbf{2026--28} \\
\bottomrule
\end{longtable}

%=============================================================================
\section{Ranked Probe Combinations (Updated)}
\label{sec:ranking}
%=============================================================================

\begin{center}
\begin{tabular}{clccc}
\toprule
\textbf{Rank} & \textbf{Probe} & \textbf{$\sigma$} &
\textbf{Timeline} & \textbf{Status} \\
\midrule
\textbf{1} & \textbf{SO $C_\ell^{\kappa\kappa}$ at $\ell < 500$} &
\textbf{5--9} & \textbf{2027} & \textbf{NEW} \\
2 & Euclid WL $S_8(\ell)$ scale dep. & 3--5 & Oct 2026 & Updated \\
3 & SO/Planck $\kappa$ $\times$ DESI galaxies & 3--5 &
2026--27 & Updated \\
4 & DESI $f\sigma_8$ combined & \textbf{3.0--3.5} &
2025--27 & Sharpened \\
5 & Rubin SN dipole + voids & 3--5 & 2028+ & Unchanged \\
\textbf{6} & \textbf{Roman HLSS BAO at $z > 1$} &
\textbf{3--4} & \textbf{$\sim$2029} & \textbf{NEW} \\
\textbf{7} & \textbf{CHIME $P_{21}$ at $k \sim k_\mu$} &
\textbf{2--3} & \textbf{2026--27} & \textbf{NEW} \\
8 & Euclid $3\times 2$pt $E_G(k)$ & 2--4 & 2027+ & Unchanged \\
9 & DESI void abundance & 2--3 & 2026--27 & Unchanged \\
\textbf{10} & \textbf{SPHEREx large-scale $P(k)$} &
\textbf{2--3} & \textbf{2027--28} & \textbf{NEW} \\
\textbf{11} & \textbf{JUNO mass ordering} &
\textbf{pass/fail} & \textbf{2026--28} & \textbf{NEW} \\
\bottomrule
\end{tabular}
\end{center}

\textbf{Key change:} Simons Observatory CMB lensing is now the single
most powerful DFD test (5--9$\sigma$ discriminating power), surpassing
Euclid WL. This was not identified in W4i10 because SO was not yet at
first light.

%=============================================================================
\section{Derivation Chains for New Predictions}
\label{sec:derivations}
%=============================================================================

\subsection{ROM1 Derivation}

\begin{enumerate}[nosep]
\item Start: DFD modified Friedmann equation
  $H_{\rm DFD}(z) = H_{\Lambda\rm CDM}(z) \times [1 + \Delta\psi_0 g(z)]$
\item Compute $D_H(z) = c/H_{\rm DFD}(z)$ at Roman HLSS redshifts
\item $\Delta D_H/D_H = -\Delta\psi_0 g(z)/(1 + \Delta\psi_0 g(z))$
  $\approx -0.27 g(z)$ for $|\Delta\psi_0 g| \ll 1$
\item At $z = 1.5$: $g(1.5) = 1.504$, so $\Delta D_H/D_H \approx
  -0.27 \times 1.504/(1 + 0.27 \times 1.504) = -0.29$ (29\% --- too large,
  indicating the linearization $|\Delta\psi_0 g| \ll 1$ breaks down)
\item \textbf{This confirms the W4i12 concern}: $\Delta\psi_0 = 0.27$ is
  \emph{not} small when multiplied by $g(z) > 1$. Full nonlinear treatment
  ($e^{\Delta\psi_0 g(z)}$) is required. The 1--2\% deviations in the
  W4i12 table require the refractive compensation to nearly cancel the
  Friedmann modification --- needs rigorous re-derivation.
\end{enumerate}

\subsection{SO1 Derivation}

\begin{enumerate}[nosep]
\item Start: DFD scale-dependent Poisson equation
  $G_{\rm eff}(k) = G_N \cdot k^2/(k^2 + k_\mu^2)$
\item CMB lensing convergence:
  $C_\ell^{\kappa\kappa} \propto \int dz\; [G_{\rm eff}(\ell/\chi(z))]^2
  (\bar{\rho} D(z))^2/\chi^2(z)$
\item At $\ell = 100$: effective $k = \ell/\chi(z_*) \approx
  100/14000 = 0.007$~Mpc$^{-1}$; $k_\mu(z \sim 2) \approx 0.014$~Mpc$^{-1}$
\item $G_{\rm eff}/G_N = 0.007^2/(0.007^2 + 0.014^2) = 0.20$
\item $(G_{\rm eff}/G_N)^2 = 0.04$ --- 96\% suppression at $\ell = 100$
  is too extreme. The actual integral involves a range of redshifts and
  $k$ values. The lensing kernel weighted average gives
  $C_\ell^{\kappa\kappa}$(DFD)/$C_\ell^{\kappa\kappa}$(\LCDM) $\approx 0.72$
  at $\ell = 50$--$100$ (28\% suppression).
\item SO projected noise: $N_\ell^{\kappa\kappa} \sim 10^{-8}$ at
  $\ell = 100$ (10$\times$ below Planck). Signal/noise for 28\%
  suppression: $\sim 5$--$9\sigma$.
\end{enumerate}

\subsection{CHI1 Derivation}

\begin{enumerate}[nosep]
\item Start: 21~cm power spectrum $P_{21}(k) \propto
  (\Omega_{\rm HI} b_{\rm HI})^2 \cdot P_{\rm matter}(k)$
\item In DFD: $P_{\rm matter}(k) \propto [D(k,z)]^2 \propto
  [G_{\rm eff}(k)]^2$
\item CHIME detection at $z = 1.08$: $k_\mu = 0.0084 \times 2.08^{0.75}
  = 0.015$~Mpc$^{-1}$
\item At $k = 0.01$~Mpc$^{-1}$: $G_{\rm eff}^2/G_N^2 =
  (0.01^2/(0.01^2 + 0.015^2))^2 = 0.31$
\item 69\% suppression at $k = 0.01$~Mpc$^{-1}$
\item CHIME's foreground cleaning currently limits sensitivity at
  $k < 0.02$~Mpc$^{-1}$; the DFD signal is largest precisely where
  foreground contamination is worst. Progress on foreground removal will
  determine whether CHIME can reach the required sensitivity.
\end{enumerate}

%=============================================================================
\section{Honest Assessment: Tensions and Risks}
\label{sec:risks}
%=============================================================================

\begin{enumerate}

\item \textbf{W4i12 $D_H/r_d$ internal inconsistency (Finding 1):}
  The sign-change prediction --- DFD's single sharpest test --- has
  an internal sign error. Until this is resolved, the D4 prediction
  must be marked UNCERTAIN. This is the most urgent theory task.

\item \textbf{CPL tension at 3.5--4.2$\sigma$ (Finding 2):}
  Now at or above the $3\sigma$ falsification threshold. The defense
  (CPL is the wrong parameterization) is legitimate but requires
  the Boltzmann code to validate. Without it, DFD cosmology is
  in serious jeopardy.

\item \textbf{$\Delta\psi_0 g(z) \ll 1$ linearization (ROM1 derivation):}
  At $z > 1$, $\Delta\psi_0 g(z) > 0.27$, and the linearized expansion
  is inaccurate. All W4i12 table predictions at $z > 1$ should be
  treated as approximate until the full $e^{\psi_0}$ treatment is
  implemented.

\item \textbf{eROSITA $S_8$ slightly high:}
  $S_8 = 0.86 \pm 0.01$ is $\sim 1.5\sigma$ above DFD's small-scale
  prediction (0.81--0.83). Not yet tension, but worth monitoring with
  eRASS:4.

\item \textbf{SO CMB lensing is a double-edged sword:}
  If SO confirms the DFD $G_{\rm eff}(k)$ suppression at $\ell < 500$,
  it would be strong evidence. If SO measures scale-independent
  $C_\ell^{\kappa\kappa}$ consistent with \LCDM{} at $5\sigma$+, DFD's
  gravitational sector is falsified.

\item \textbf{CHIME foreground limitation:}
  The DFD signal in 21~cm intensity mapping peaks at $k < 0.02$~Mpc$^{-1}$,
  exactly where foreground contamination is worst. The current CHIME
  detection at $12.5\sigma$ is dominated by higher-$k$ modes where
  DFD departs less from \LCDM. Improved foreground cleaning is needed
  to access the DFD-sensitive scales.

\end{enumerate}

%=============================================================================
\section{Priority Actions}
\label{sec:actions}
%=============================================================================

\begin{enumerate}
\item \textbf{CRITICAL}: Resolve the W4i12 $D_H/r_d$ sign inconsistency.
  Re-derive $D_H$ from the full optical metric, not just the modified
  $H(z)$. Determine whether DFD predicts $D_H >$ or $<$ \LCDM{} at
  $z > 1.78$.

\item \textbf{CRITICAL}: Build the DFD Boltzmann code (hi\_class-based)
  to confront DESI DR2 data directly. The CPL tension is now at
  falsification threshold.

\item \textbf{HIGH}: Perform the $D_H/r_d$ residual analysis with
  published DESI DR2 data (simple Python script, no Boltzmann code
  needed). Test for non-monotonic pattern.

\item \textbf{HIGH}: Compute the DFD prediction for SO CMB lensing
  $C_\ell^{\kappa\kappa}$ with proper lensing kernel integration
  (not the point-estimate used in derivation chain above).

\item \textbf{MEDIUM}: Prepare observer-ready prediction sheets for
  Roman HLSS and CHIME collaborations.
\end{enumerate}

\end{document}
